\documentclass[12pt,reqno]{amsart}
\usepackage{math-review}
\providecommand{\Ric}{\operatorname{Ric}}
\providecommand{\Hess}{\operatorname{Hess}}
\providecommand{\tr}{\operatorname{tr}}
\providecommand{\Sym}{\operatorname{Sym}}

\title{From a uniform H\"older bound to a valid limiting equation}
\date{}
\begin{document}
\maketitle

An a priori estimate is useful only after its topology and constants are
specified. This lecture proves the compactness step and then separates
the two further tasks: preserving the equation's domain and passing its residual.

\section{The conditional target}
We assume multivariable calculus, compactness of bounded real sequences
and smooth partitions of unity. The compactness, interpolation and
admissibility arguments below are developed explicitly.

\begin{definition}\label{def:norm}
Fix a compact smooth $d$-manifold $M$ without boundary, a smooth background metric,
a finite smooth atlas $(U_a,\phi_a)$ and a smooth partition of unity $(\chi_a)$
with $\operatorname{supp}\chi_a\Subset U_a$.
For $0<\gamma<1$, extend $v_a=(\chi_a u)\circ\phi_a^{-1}$ by zero to $\mathbb R^d$ and set
\[
 \|u\|_{2,\gamma}=\sum_a\left(\sum_{|\nu|\leq2}\|D^\nu v_a\|_\infty
       +\sum_{|\nu|=2}[D^\nu v_a]_\gamma\right),
 \qquad [w]_\gamma=\sup_{x\ne y}\frac{|w(x)-w(y)|}{|x-y|^\gamma}.
\]
Here $\nu$ is a multi-index, and $C^{2,\gamma}(M)$ consists of the $C^2$ functions
for which this norm is finite. All these choices stay fixed throughout a sequence.
\end{definition}
\begin{theorem}\label{thm:closedness}
Fix the compact smooth manifold and norms just defined, integers $d,r\geq1$,
and $0<\beta<\alpha<1$. Let
\[
 A:C^{2,\beta}(M)\longrightarrow C^0(M;\Sym_r(\mathbb R))
\]
be continuous for the uniform operator norm in the target. Define
\[
 \mathcal U=\{u\in C^{2,\beta}(M):A(u)(x)>0\ \text{for every }x\in M\},
 \qquad \mathcal F:[0,1]\times\mathcal U\longrightarrow C^0(M),
\]
and assume $\mathcal F$ is jointly continuous. Suppose $t_j\to t_*$ and
$u_j\in C^{2,\alpha}(M)\cap\mathcal U$ satisfy, with fixed $C<\infty$ and $\delta>0$,
\[
 \mathcal F(t_j,u_j)=0,\qquad \|u_j\|_{2,\alpha}\leq C,\qquad
 A(u_j)(x)\geq\delta I_r\quad\text{for every }j,x.
\]
Then some subsequence converges in $C^{2,\beta}(M)$ to $u_*\in C^{2,\alpha}(M)$, with
\[
 \|u_*\|_{2,\alpha}\leq C,\quad A(u_*)(x)\geq\delta I_r\quad(x\in M),
 \quad u_*\in\mathcal U,\quad\mathcal F(t_*,u_*)=0.
\]
This is conditional on the displayed bounds and continuity; it asserts neither
existence of such solutions nor higher regularity or convergence in $C^{2,\alpha}$.
\end{theorem}

The estimate and the positive constant $\delta$ are assumptions.
A specific PDE must supply them, uniformly in $j$, from its own
coefficients, geometry and normalization; this theorem does not do so.

The proof proceeds from fixed coordinate pieces to a single limit,
then to $C^{2,\beta}$ convergence. Only after the margin has placed the
limit in $\mathcal U$ do we invoke continuity of $\mathcal F$.

\section{Constructing the convergent subsequence}
Write $v_{a,j}=(\chi_a u_j)\circ\phi_a^{-1}$, extended by zero.
The chart supports are fixed compact sets. The assumed norm bound gives
\[
 \|D^\nu v_{a,j}\|_\infty\leq C\quad(|\nu|\leq2),
 \qquad [D^\nu v_{a,j}]_\alpha\leq C\quad(|\nu|=2).
\]
The second derivatives have a common H\"older modulus. Lower derivatives
have a common Lipschitz modulus by the next-derivative bound and
integration along a segment in $\mathbb R^d$.

\subsection*{A common subsequence}
For each of the finitely many derivative families, choose a countable
dense set in its compact support. Enumerate all scalar evaluations and
successively extract subsequences on which the next evaluation converges.

Take the diagonal of these nested subsequences. It converges at every
listed point. To promote this to uniform convergence, choose a finite
net from the dense set using the common modulus of continuity.

For any two late terms $w_j,w_k$ of one derivative family and a net
point $s$ close to $x$, use
\[
 |w_j(x)-w_k(x)|\leq
 |w_j(x)-w_j(s)|+|w_j(s)-w_k(s)|+|w_k(s)-w_k(x)|.
\]
The first and third terms are small by the common modulus. The middle
term is uniformly small on the finite net, so the subsequence is uniformly
Cauchy. This is the scalar Arzel\`a--Ascoli argument of \cite{S-Shapiro}.

\subsection*{The limits really are derivatives}
Uniform limits of separately named derivatives need compatibility.
For $|\nu|\leq1$, keep the identity
\[
 D^\nu v_{a,j}(x+he_i)-D^\nu v_{a,j}(x)
 =\int_0^h D^{\nu+e_i}v_{a,j}(x+se_i)\,ds.
\]
Uniform convergence on the segment permits passage under the integral.
The limiting identity identifies first and second derivatives, so each
local limit is $C^2$ with the derivatives just extracted.

Fixed smooth transition rules pass to the limit as well. The partition
of unity glues these compatible pieces to a global function $u_*$.
Compact support inside each chart makes zero extension harmless.

For any fixed pair $x\ne y$, passing to the limit in its H\"older quotient
retains the $\alpha$ bound. Taking suprema and then the finite sum yields
\[
 u_*\in C^{2,\alpha}(M),\qquad \|u_*\|_{2,\alpha}\leq C.
\]

\section{Why the exponent must be weakened}
Uniform convergence of second derivatives does not by itself control
their $\beta$ seminorms. For a bounded $\alpha$-H\"older function $w$,
two available bounds can be combined:
\[
 |w(x)-w(y)|\leq
 \min\{2\|w\|_\infty,[w]_\alpha|x-y|^\alpha\}.
\]
For positive $a,b$ and $0<\theta<1$, $\min(a,b)\leq a^{1-\theta}b^\theta$.
Set $\theta=\beta/\alpha$ and divide by $|x-y|^\beta$:
\begin{equation}\label{eq:interpolation}
 [w]_\beta\leq(2\|w\|_\infty)^{1-\beta/\alpha}
                  [w]_\alpha^{\beta/\alpha}.
\end{equation}
The zero cases follow directly. Apply this to each difference
$D^\nu(v_{a,j}-v_{a,*})$ with $|\nu|=2$: its sup norm tends to zero
and its $\alpha$ seminorm is at most $2C$.

The strictly positive exponent $1-\beta/\alpha$ forces the right side
of \eqref{eq:interpolation} to vanish. Combined with the lower derivatives,
this proves
\[
 \|u_j-u_*\|_{2,\beta}\longrightarrow0
\]
along the extracted subsequence. The theorem's stronger regularity of
the limit is a boundedness conclusion, not convergence in that stronger norm.

\section{Passing the domain before the equation}
Continuity of $A$ now supplies uniform matrix convergence. For every
$x\in M$ and $v\in\mathbb R^r$,
\[
 v^TA(u_*)(x)v=\lim_j v^TA(u_j)(x)v\geq\delta|v|^2.
\]
Thus $A(u_*)\geq\delta I_r>0$, and $u_*\in\mathcal U$. The pair
$(t_*,u_*)$ is therefore in the actual domain of $\mathcal F$, where
joint continuity gives
\[
 \mathcal F(t_*,u_*)=\lim_j\mathcal F(t_j,u_j)=0\quad\text{in }C^0(M).
\]
This completes Theorem~\ref{thm:closedness}. If the sequence also has
mean zero, that normalization passes because
\[
 \left|\int_M(u_j-u_*)\,dV_g\right|
 \leq\operatorname{Vol}_g(M)\|u_j-u_*\|_\infty\longrightarrow0.
\]

\subsection*{A concrete continuity check}
On a flat torus take a fixed continuous $h$ and
\[
 \mathcal F(t,u)=\operatorname{tr}(D^2u)+u^3-th.
\]
The second derivatives converge uniformly. On the common bounded range
$|u_j|,|u_*|\leq R$, the nonlinear term satisfies
\[
 \|u_j^3-u_*^3\|_\infty\leq3R^2\|u_j-u_*\|_\infty.
\]
The parameter term is bounded by $|t_j-t_*|\|h\|_\infty$.
These three estimates verify the claimed $C^0$ continuity for this
example; they provide no existence theorem or a priori solution bound.

\section{Localization and failure models}
On an open subset of $\mathbb R^d$, estimates on every relatively compact
open set give a local conclusion. Choose an exhaustion $K_\ell$ with
$K_\ell\Subset\operatorname{int}K_{\ell+1}$ and insert fixed cutoffs.

Extract on $K_1$, then on $K_2$ inside that subsequence, and continue.
The diagonal converges on every fixed $K_\ell$; limits agree on overlaps.
The bounds may depend on $\ell$, so the conclusion is only $C^{2,\beta}_{\rm loc}$.

For a nonzero compactly supported smooth function $\eta$ on $\mathbb R$,
$u_j(x)=\eta(x-j)$ eventually vanishes on every fixed compact set, but
\[
 \|u_j\|_\infty=\|\eta\|_\infty>0.
\]
Thus local extraction does not imply global uniform convergence.
It also supplies no boundary convergence on a noncompact or open domain.

\subsection*{Losing the positive margin}
Let $A(u)=uI_r$, let $\mathcal U$ be the positive functions, and set
$\mathcal F(t,u)=u-t$. The constant functions
\[
 t_j=u_j=1/j
\]
have uniform H\"older bounds and solve the equation, but their limit
is zero and leaves $\mathcal U$. Individual strict inequalities cannot
replace a uniform $\delta>0$.

\subsection*{Losing the exponent gap}
On $\mathbb R/(2\pi\mathbb Z)$, let $u_n(x)=n^{-2-\alpha}\cos(nx)$.
Its second derivative is $-n^{-\alpha}\cos(nx)$ and converges uniformly
to zero. The bound $|\cos a-\cos b|\leq\min\{2,|a-b|\}$ gives uniform
$\alpha$ seminorms.

At $x=0$ and $y=\pi/n$, however, the second-derivative quotient is
\[
 \frac{|u_n''(0)-u_n''(\pi/n)|}{(\pi/n)^\alpha}
       =\frac{2}{\pi^\alpha}.
\]
It cannot converge to zero in $C^{2,\alpha}$. Since every possible
uniform limit is zero, no subsequence can repair that failure.

\subsection*{Weak convergence is a different input}
Work in the real space $L^2(0,2\pi;\mathbb R)$, with inner product
$\langle f,g\rangle=\int_0^{2\pi}fg\,dx$ and norm
$\|g\|_2^2=\langle g,g\rangle$.

By definition, $f_n\rightharpoonup0$ weakly in this space means
$\int_0^{2\pi}f_ng\,dx\to0$ for every $g\in L^2(0,2\pi;\mathbb R)$.

Set $e_n(x)=\sin(nx)/\sqrt\pi$ for integers $n\geq1$.
Integrating the product-to-sum identity
\[
 2\sin(nx)\sin(mx)=\cos((n-m)x)-\cos((n+m)x)
\]
gives $\langle e_n,e_m\rangle=0$ for $n\ne m$ and $1$ for $n=m$.
Thus the sequence is orthonormal.

Fix any $g$ in this real $L^2$ space and put $a_n=\langle g,e_n\rangle$.
Expanding a finite squared norm using orthonormality gives
\[
 0\leq\left\|g-\sum_{n=1}^N a_ne_n\right\|_2^2
       =\|g\|_2^2-\sum_{n=1}^N|a_n|^2.
\]
Hence the partial sums of $\sum|a_n|^2$ are bounded, so $a_n\to0$.
For this arbitrary $g$, $\int_0^{2\pi}\sin(nx)g(x)\,dx=\sqrt\pi\,a_n\to0$,
proving $\sin(nx)\rightharpoonup0$.

The squares also lie in $L^2$. Testing against $1\in L^2$ gives
\[
 \int_0^{2\pi}\sin^2(nx)\,dx=\pi.
\]
Thus the squares do not converge weakly to zero. The nonlinear passage
in the theorem instead uses strong convergence in a topology for which
continuity was actually verified.

\bibliographystyle{../manuscript/localrefs}
\bibliography{../generated/references}
\end{document}
